\documentclass[border=10pt]{standalone}
\usepackage{fontspec}
\setmainfont{Noto Serif}[Ligatures=TeXOff]
\XeTeXgenerateactualtext=1
\newfontfamily\bengalifont{Noto Serif Bengali}[Script=Bengali, NFSSFamily=bnfam]
\usepackage[Latin,Bengali]{ucharclasses}
% group-free font switching (as in header-common.tex): robust inside braces, cells and quotes
\makeatletter
\def\BN@saved{\rmdefault}
\def\BN@bn{bnfam}
\newcommand\BNin{\edef\BN@tmp{\f@family}\ifx\BN@tmp\BN@bn\else\let\BN@saved\BN@tmp\fi\fontfamily{bnfam}\selectfont}
\newcommand\BNout{\fontfamily{\BN@saved}\selectfont}
\makeatother
\setTransitionsFor{Bengali}{\BNin}{\BNout}
% Typographic quotes and dashes are not in a ucharclasses class (class 0), so a Bangla word right
% after “ or ‘ would not trigger the Bangla font. Treat them as Latin so that every switch fires.
\XeTeXcharclass"2018=\XeTeXcharclass`A
\XeTeXcharclass"2019=\XeTeXcharclass`A
\XeTeXcharclass"201C=\XeTeXcharclass`A
\XeTeXcharclass"201D=\XeTeXcharclass`A
\XeTeXcharclass"2013=\XeTeXcharclass`A
\XeTeXcharclass"2014=\XeTeXcharclass`A
\XeTeXcharclass"2026=\XeTeXcharclass`A
\usepackage{newunicodechar}
\newunicodechar{।}{{\bengalifont\char"0964}}
\newunicodechar{॥}{{\bengalifont\char"0965}}

\usepackage[table]{xcolor}
\usepackage{booktabs,tabularx,array}
\definecolor{bapilrule}{HTML}{1F2937}
\definecolor{bapilstripe}{HTML}{F3F4F6}
\renewcommand{\arraystretch}{1.28}
\setlength{\tabcolsep}{5pt}

\begin{document}
\begin{minipage}{22cm}
{\large\bfseries সারণি ৬৪।}\ {\large শাসন-সংস্থা: অবস্থা, নিয়োগ, সিদ্ধান্তের অধিকার ও সুরক্ষা}\par\vspace{7pt}
\small
\rowcolors{2}{bapilstripe}{white}
\begin{tabularx}{\linewidth}{@{}>{\hsize=1.031\hsize\raggedright\arraybackslash}X>{\hsize=0.977\hsize\raggedright\arraybackslash}X>{\hsize=1.194\hsize\raggedright\arraybackslash}X>{\hsize=0.705\hsize\raggedright\arraybackslash}X>{\hsize=1.194\hsize\raggedright\arraybackslash}X>{\hsize=0.868\hsize\raggedright\arraybackslash}X>{\hsize=1.031\hsize\raggedright\arraybackslash}X@{}}
\arrayrulecolor{bapilrule}\toprule
\rowcolor{white}\textbf{সংস্থা} & \textbf{অবস্থা ও ভিত্তি} & \textbf{সদস্যপদ ও নিয়োগ} & \textbf{সভাপতি ও সচিবালয়} & \textbf{সিদ্ধান্তের অধিকার ও সীমা} & \textbf{প্রতিবেদন ও স্বচ্ছতা} & \textbf{স্বাধীনতা ও স্বার্থের দ্বন্দ্ব} \\
\midrule
\textbf{স্টিয়ারিং বোর্ড} & প্রস্তাবিত; সমঝোতা স্মারক ও স্বাগতিক-চুক্তির মাধ্যমে গঠিত; কোনো সংবিধিবদ্ধ ক্ষমতা নেই & নিজ নিজ প্রতিষ্ঠানের মনোনীত ৭ থেকে ৯ জন সদস্য; একাডেমিক ও নাগরিক সমাজের সদস্য উন্মুক্ত আহ্বানে নির্বাচিত; মেয়াদ ৩ বছর & আইসিটি বিভাগ ও স্বাগতিক প্রতিষ্ঠান থেকে যুগ্ম সভাপতি; পরিচালকের দপ্তরে সচিবালয় & পরিকল্পনা, বাজেট, অংশীদারত্ব ও জ্যেষ্ঠ নিয়োগ অনুমোদন; সংস্থার সঙ্গে যৌথ G4 সিদ্ধান্ত; ERRC-এর প্রত্যাখ্যান বা আইন অগ্রাহ্য করতে পারে না & ত্রৈমাসিক সভা; সিদ্ধান্ত ও কার্যবিবরণী সংক্ষেপে প্রকাশিত; বার্ষিক প্রতিবেদন & বার্ষিক ঘোষণা; সংশ্লিষ্ট সিদ্ধান্ত থেকে বিরত থাকা; শিল্পখাতের সদস্য নিজ খাতের ক্রয়ে ভোট দেন না \\
\textbf{নৈতিকতা ও ঝুঁকি পর্যালোচনা কমিটি (ERRC)} & প্রস্তাবিত; সনদ দ্বারা গঠিত; সনদের মাধ্যমে BAPIL-এর ওপর বাধ্যতামূলক, আর স্বাগতিকের ওপর কেবল তার সংবিধি ও স্বাগতিক-চুক্তি যতটুকু অনুমোদন করে; আইন দ্বারা নয় & ৫ থেকে ৭ জন সদস্য, অধিকাংশ বহিঃস্থ; স্বাধীন মনোনয়ন প্যানেলের উন্মুক্ত আহ্বানে নির্বাচিত ও স্টিয়ারিং বোর্ড কর্তৃক নিশ্চিত; মেয়াদ ৩ বছর, একবার নবায়নযোগ্য; শুনানির পর বোর্ডের দুই-তৃতীয়াংশের সিদ্ধান্তে কেবল যুক্তিসংগত কারণে অপসারণ & সদস্যদের নির্বাচিত বহিঃস্থ সভাপতি; শাসন-বিশ্লেষক সচিবালয়ের দায়িত্বে & G1, G2 ও G4-এ T2 ও T3 স্তরের কাজ অনুমোদন, শর্তসাপেক্ষ অনুমোদন, প্রত্যাখ্যান বা স্থগিত; প্রত্যাখ্যান বাধ্যতামূলক; আইন বা স্বাগতিক যা নিষেধ করে তা অনুমোদন করতে পারে না & ৩০ দিনের মধ্যে সিদ্ধান্তের সারাংশ; ভিন্নমত লিপিবদ্ধ; বার্ষিক প্রতিবেদন & অধিকাংশ বহিঃস্থ; মূল্যায়নাধীন সরবরাহকারীর কোনো কর্মী সদস্য নন; সুরক্ষিত বাজেট; কর্মীরা সরাসরি উদ্বেগ জানাতে পারেন \\
\textbf{কারিগরি স্থাপত্য বোর্ড (TAB)} & প্রস্তাবিত; পরিচালকের অধীন অভ্যন্তরীণ কারিগরি সংস্থা & প্রধান প্রকৌশলীবৃন্দ, নিরাপত্তা প্রধান এবং বোর্ডের অনুমোদনে পরিচালক কর্তৃক নিযুক্ত দুজন বহিঃস্থ পর্যালোচক & প্রধান প্রকৌশলী; প্ল্যাটফর্ম দলে সচিবালয় & G1-এ নকশা, মডেল ও প্রধান নির্ভরতা অনুমোদন; ADR রক্ষণাবেক্ষণ & রিপোজিটরিতে ADR উন্মুক্ত; মাসিক সভা & বহিঃস্থ পর্যালোচকেরা বাণিজ্যিক স্বার্থ ঘোষণা করেন এবং নিজ নিয়োগকর্তার পণ্য পর্যালোচনা করেন না \\
\textbf{আন্তর্জাতিক উপদেষ্টা প্যানেল} & প্রস্তাবিত; শুধু পরামর্শমূলক & স্টিয়ারিং বোর্ডের আমন্ত্রণে; ব্যক্তিগত সক্ষমতায়; মেয়াদ ২ বছর & সদস্যদের বাছাই করা সভাপতি; পরিচালকের দপ্তরে সচিবালয় & কোনো সিদ্ধান্ত-ক্ষমতা নেই; পদ্ধতি, বেঞ্চমার্ক ও শাসন বিষয়ে পরামর্শ দেয় & বার্ষিক নোট প্রকাশিত & স্বার্থ ঘোষণা \\
\textbf{পরিচালক} & প্রস্তাবিত পদ; স্বাগতিক প্রতিষ্ঠানের কর্মী & উন্মুক্ত প্রতিযোগিতার মাধ্যমে স্বাগতিক কর্তৃক নিযুক্ত ও স্টিয়ারিং বোর্ড কর্তৃক নিশ্চিত & প্রযোজ্য নয় & অনুমোদিত পরিকল্পনা ও বাজেটের মধ্যে পরিচালনাগত সিদ্ধান্ত; যেকোনো পাইলট বন্ধ করতে পারেন & বোর্ডের কাছে মাসিক প্রতিবেদন & স্বাগতিকের চাকরিবিধি; ঘোষণা \\
\textbf{স্বাগতিক সংস্থার প্রোডাক্ট ওনার} & নিজ সংস্থা কর্তৃক লিখিতভাবে মনোনীত কর্মকর্তা & সংস্থাপ্রধান কর্তৃক মনোনীত & প্রযোজ্য নয় & চাহিদা নির্ধারণ; নিজ সংস্থায় G3 ও G4-এ পাইলট গ্রহণ বা প্রত্যাখ্যান & পাইলট প্রতিবেদন & সংস্থার বিধি অনুযায়ী কাজ করেন; BAPIL থেকে কোনো পারিশ্রমিক নেন না \\
\textbf{স্বাধীন মনোনয়ন প্যানেল} & প্রস্তাবিত; প্রতিটি ERRC মেয়াদের জন্য একবার আহূত & আইন, নৈতিকতা ও নাগরিক সমাজে স্বীকৃত তিনজন ব্যক্তি, স্বাগতিক প্রতিষ্ঠানের গবেষণা-নৈতিকতা কমিটির আমন্ত্রণে & সদস্যরা একজন আহ্বায়ক বাছাই করেন & ERRC সদস্যদের সংক্ষিপ্ত তালিকা ও সুপারিশ & আহ্বান ও নির্বাচনের মানদণ্ড প্রকাশ & সদস্যরা নিজেদের বাছাই করা ERRC-তে কাজ করতে পারেন না \\
\bottomrule
\end{tabularx}
\par\vspace{5pt}{\footnotesize\itshape এসব সংস্থার কোনোটিই এখনো গঠিত হয়নি। BAPIL-এর অভ্যন্তরীণ নিয়ম আইনের পাশাপাশি প্রযোজ্য, কখনো আইনের বিকল্প নয়; আইন, কোনো নিয়ন্ত্রক বা স্বাগতিক প্রতিষ্ঠান ভিন্ন কিছু চাইলে আইনগত প্রয়োজনীয়তাই প্রাধান্য পাবে।}
\par\vspace{3pt}{\scriptsize\color{gray} উৎস: BAPIL টেকনিক্যাল ব্লুপ্রিন্ট, সংস্করণ ১.০ (২০২৬)। CC BY 4.0।}
\end{minipage}
\end{document}
